%!TEX TS-program = xelatex
%!TEX encoding = UTF-8 Unicode
% Deedy CV/Resume Template
% Original author: Debarghya Das (http://debarghyadas.com)
% Repository: https://github.com/deedy/Deedy-Resume
% License: Apache 2.0
%
% IMPORTANT: Compile with XeLaTeX (requires fontspec + Lato/Raleway fonts)
%

\documentclass[]{deedy-resume-openfont}
\usepackage{fancyhdr}
\usepackage{fontawesome5}
\usepackage{pifont}
\newcommand{\githubstars}[1]{[{\faGithubSquare}\,\ding{72}\,\textbf{#1}]}

\pagestyle{fancy}
\fancyhf{}

\begin{document}

\lastupdated

%%%%%%%%%%%%%%%%%%%%%%%%%%%%%%%%%%%%%%
%
%     TITLE NAME
%
%%%%%%%%%%%%%%%%%%%%%%%%%%%%%%%%%%%%%%
\namesection{Hyungtae}{Lim}{
  \urlstyle{same}\href{https://limhyungtae.github.io/hyungtae-lim/}{limhyungtae.github.io} |
  \href{mailto:shapelim@mit.edu}{shapelim@mit.edu} |
  \href{https://scholar.google.com/citations?user=S1A3nbIAAAAJ}{Google Scholar}
}

%%%%%%%%%%%%%%%%%%%%%%%%%%%%%%%%%%%%%%
%
%     COLUMN ONE (LEFT)
%
%%%%%%%%%%%%%%%%%%%%%%%%%%%%%%%%%%%%%%
\begin{minipage}[t]{0.33\textwidth}

%%%%%%%%%%%%%%%%%%%%%%%%%%%%%%%%%%%%%%
%     EDUCATION
%%%%%%%%%%%%%%%%%%%%%%%%%%%%%%%%%%%%%%
\section{Education}

\subsection{KAIST}
\descript{Ph.D. in EE \& Robotics}
\location{Mar. 2020 -- Feb. 2023 | Daejeon, Korea}
Advisor: Prof. Hyun Myung \\
Dissertation: Robust LiDAR SLAM \\
for Autonomous Vehicles
\sectionsep

\subsection{KAIST}
\descript{M.S. in EE \& Robotics}
\location{Mar. 2018 -- Feb. 2020 | Daejeon, Korea}
\sectionsep

\subsection{KAIST}
\descript{B.S. in Mechanical Engineering}
\location{Mar. 2013 -- Aug. 2018 | Daejeon, Korea}
Dean's List (Fall 2015)
\sectionsep

%%%%%%%%%%%%%%%%%%%%%%%%%%%%%%%%%%%%%%
%     LINKS
%%%%%%%%%%%%%%%%%%%%%%%%%%%%%%%%%%%%%%
\section{Links}
Website:// \href{https://limhyungtae.github.io/hyungtae-lim/}{\bf limhyungtae.github.io} \\
GitHub:// \href{https://github.com/limhyungtae}{\bf limhyungtae} \\
LinkedIn:// \href{https://www.linkedin.com/in/hyungtae-lim}{\bf hyungtae-lim} \\
Scholar:// \href{https://scholar.google.com/citations?user=S1A3nbIAAAAJ}{\bf Hyungtae Lim}
\sectionsep

%%%%%%%%%%%%%%%%%%%%%%%%%%%%%%%%%%%%%%
%     SKILLS
%%%%%%%%%%%%%%%%%%%%%%%%%%%%%%%%%%%%%%
\section{Skills}
\subsection{Programming}
\location{Proficient:}
C++ \textbullet{} Python \textbullet{} ROS/ROS2 \\
\location{Experienced:}
PyTorch \textbullet{} CUDA \textbullet{} CMake \\
OpenCV \textbullet{} \LaTeX
\sectionsep

\subsection{Research}
SLAM \textbullet{} 3D LiDAR \\
Point Cloud Registration \\
Ground Segmentation \\
Sensor Fusion \textbullet{} Deep Learning \\
Multi-Robot Mapping
\sectionsep

\subsection{Languages}
Korean (Native) \textbullet{} English (Fluent)
\sectionsep

\end{minipage}
\hfill
%%%%%%%%%%%%%%%%%%%%%%%%%%%%%%%%%%%%%%
%
%     COLUMN TWO (RIGHT)
%
%%%%%%%%%%%%%%%%%%%%%%%%%%%%%%%%%%%%%%
\begin{minipage}[t]{0.66\textwidth}

%%%%%%%%%%%%%%%%%%%%%%%%%%%%%%%%%%%%%%
%     EXPERIENCE
%%%%%%%%%%%%%%%%%%%%%%%%%%%%%%%%%%%%%%
\section{Experience}

\runsubsection{SPARK Lab, MIT}
\descript{| Postdoctoral Associate}
\location{Apr. 2024 -- Mar. 2026 | Cambridge, MA}
\vspace{0.5\topsep}
\begin{tightemize}
\item Lead research on large-scale multi-robot SLAM enabling real-time, consistent 3D map building across robots and sessions in complex urban environments.
\item Mentor researchers on 3D perception, autonomous mapping, and scene understanding; maintainer of 4 open-source robotics repos (\githubstars{2.1K}--\githubstars{643}).
\end{tightemize}
\sectionsep

\runsubsection{Urban Robotics Lab, KAIST}
\descript{| Postdoctoral Fellow}
\location{Mar. 2023 -- Mar. 2024 | Daejeon, Korea}
\begin{tightemize}
\item Developed real-time LiDAR-inertial and visual SLAM pipelines for robust autonomous navigation in dynamic real-world environments.
\item Designed deep learning perception modules for dynamic obstacle detection, improving mapping reliability in crowded and unstructured scenes.
\end{tightemize}
\sectionsep

\runsubsection{StachnissLab, Univ. Bonn}
\descript{| Visiting Scholar}
\location{Nov. 2022 -- Feb. 2023 | Bonn, Germany}
\begin{tightemize}
\item Built ERASOR2: instance-aware dynamic object removal enabling high-fidelity static LiDAR map building in dynamic urban environments (RSS 2023).
\end{tightemize}
\sectionsep

\runsubsection{NAVER LABS}
\descript{| Research Intern}
\location{Apr. 2022 -- Sep. 2022 | Gyeonggi-Do, Korea}
\begin{tightemize}
\item Engineered ORORA, a radar odometry system robust to outliers enabling reliable localization in GPS-denied and featureless environments (IEEE ICRA 2023).
\end{tightemize}
\sectionsep

%%%%%%%%%%%%%%%%%%%%%%%%%%%%%%%%%%%%%%
%     KEY PROJECTS
%%%%%%%%%%%%%%%%%%%%%%%%%%%%%%%%%%%%%%
\section{Key Open-Source Research Projects}
\vspace{0.3cm}
\begin{tightemize}
\item \href{https://github.com/MIT-SPARK/KISS-Matcher}{\textbf{KISS-Matcher}} \githubstars{643} -- Real-time 3D registration pipeline for large-scale multi-session LiDAR SLAM (\textit{IEEE ICRA 2025})
\item \href{https://github.com/MIT-SPARK/VGGT-SLAM}{\textbf{VGGT-SLAM}} \githubstars{882} -- Dense real-time RGB SLAM with globally consistent 3D scene reconstruction (\textit{NeurIPS 2025})
\item \href{https://github.com/url-kaist/patchwork-plusplus}{\textbf{Patchwork++}} \githubstars{842} -- Real-time ground segmentation deployed in AV and mobile robotics production pipelines (\textit{IROS 2022})
\item \href{https://github.com/LimHyungTae/ERASOR}{\textbf{ERASOR}} \githubstars{494} -- Dynamic object removal system for high-fidelity static LiDAR map building (\textit{IEEE RA-L 2021})
\end{tightemize}
\sectionsep

%%%%%%%%%%%%%%%%%%%%%%%%%%%%%%%%%%%%%%
%     PUBLICATIONS
%%%%%%%%%%%%%%%%%%%%%%%%%%%%%%%%%%%%%%
\section{Publications}
10\texttt{+} first-author papers: IEEE ICRA, RSS, RA-L, NeurIPS, ICCV, IJRR
\sectionsep

%%%%%%%%%%%%%%%%%%%%%%%%%%%%%%%%%%%%%%
%     AWARDS
%%%%%%%%%%%%%%%%%%%%%%%%%%%%%%%%%%%%%%
\section{Awards}
\begin{tabular}{rl}
2025 & Outstanding Reviewer Award, IEEE ICRA'25 (5/7,641) \\
2024 & RSS Pioneers 2024 (30/202 selected) \\
2024 & 1\textsuperscript{st} Prize, HILTI SLAM Challenge, IEEE ICRA'24 \\
2023 & Best Paper Award, IEEE RA-L (5 of 1,100 papers) \\
2023 & 1\textsuperscript{st} Prize, HILTI SLAM Challenge, IEEE ICRA'23 \\
2023 & CES'23 Innovation Award \\
2022 & 2\textsuperscript{nd} Prize, HILTI SLAM Challenge, IEEE ICRA'22 \\
2020 & Student Best Paper Award, ICCAS \\
\end{tabular}
\sectionsep

\end{minipage}
\end{document}
